\documentclass[letterpaper]{article}

\usepackage[margin=1.5cm]{geometry}
\usepackage{amsmath,amsfonts}
\usepackage{url}
\usepackage{bm}

\newcommand{\Real}{\mathbb{R}}
\newcommand{\tr}{^\top}
\newcommand{\opt}{^\star}

\setlength{\parskip}{3mm plus1mm minus1mm}
\setlength{\parindent}{0mm}

\title{Problem Set 8}
\author{Foundations of Machine Learning}
\date{}

\begin{document}
\maketitle

\begin{itemize}
\item \textbf{Submission}: Turn in as a \texttt{PDF} and a \textbf{separate source code file} on \url{http://mycourses.unh.edu}. Scanned hand-written notes are encouraged.
\item \textbf{Questions}: See the class website for a link to office hours
\end{itemize}

% Covers: Linear programming (CO 4.3, ILA 10.4) and Machine learning: linear regression (ISL 2--3).

\textbf{Relevant reading}: CO 4.3, ILA 10.4 (linear programming); ISL 2--3 (linear regression). Problems marked \textbf{[code]} should be submitted as a program in a language of your choice (Julia or Python recommended). The cited problem numbers are from ILA, 5th edition (Strang, 2016), and ISL, 2nd edition (James et al., 2021).

\begin{enumerate}
\item ILA, Problem Set 10.4 (linear programming): problems 1, 3.
\item \textbf{[code]} Formulate and solve a small revenue-maximization linear program (using, e.g., JuMP in Julia or PuLP/SciPy in Python). Report the optimal solution, the optimal value, and which constraints are active.
\item ISL, Chapter 3, Conceptual exercises: problems 1, 3, 4.
\item ISL, Chapter 3, Applied exercises: the simple linear regression exercise on the \texttt{Auto} data set (problem 8) \textbf{[code]}.
\item \textbf{Time spent.} Report approximately how many hours you spent completing this problem set (including the reading).
\end{enumerate}

\end{document}
